\documentclass[12pt,a4paper]{article}

% ── Packages ───────────────────────────────────────────────────────────────
\usepackage[utf8]{inputenc}
\usepackage[T1]{fontenc}
\usepackage[english]{babel}
\usepackage{geometry}
\usepackage{setspace}
\usepackage{amsmath,amssymb}
\usepackage{booktabs}
\usepackage{longtable}
\usepackage{graphicx}
\usepackage{hyperref}
\usepackage{natbib}
\usepackage{xcolor}
\usepackage{caption}
\usepackage{subcaption}
\usepackage{multirow}
\usepackage{array}
\usepackage{tabularx}
\usepackage{lmodern}
\usepackage{microtype}
\usepackage{url}
\usepackage{listings}
\usepackage{fancyhdr}
\usepackage{abstract}
\usepackage{titlesec}
\usepackage{footnote}

\geometry{
  top=2.5cm, bottom=2.5cm,
  left=3.0cm, right=3.0cm,
  headheight=15pt
}

\doublespacing

\hypersetup{
  colorlinks=true,
  linkcolor=black,
  citecolor=black,
  urlcolor=blue
}

\lstset{
  basicstyle=\ttfamily\small,
  breaklines=true,
  frame=none,
  xleftmargin=1em
}

\pagestyle{fancy}
\fancyhf{}
\rhead{\thepage}
\lhead{\small\textit{Measuring Political and Economic Uncertainty in Peru}}
\renewcommand{\headrulewidth}{0.4pt}

% ── Title ──────────────────────────────────────────────────────────────────
\title{%
  \Large\textbf{Measuring Political and Economic Uncertainty in Peru:\\
  A Daily News-Based Index}\\[1em]
  \normalsize Preliminary Draft --- March 2026\\
  \normalsize Please do not cite without permission
}

\author{%
  Carlos C\'{e}sar Ch\'{a}vez Padilla\thanks{%
    Center for the Economics of Human Development, University of Chicago.
    E-mail: \texttt{cchavezpadilla@uchicago.edu}.
    I thank the team at Qhawarina.pe for data infrastructure support.
    All errors are my own.}
}

\date{March 2026}

% ── Document ───────────────────────────────────────────────────────────────
\begin{document}

\maketitle
\thispagestyle{empty}

% ── Abstract ───────────────────────────────────────────────────────────────
\begin{abstract}
\noindent
We construct two daily indices of political and economic uncertainty for
Peru---the \emph{\'Indice de Riesgo Pol\'itico} (IRP) and the
\emph{\'Indice de Riesgo Econ\'omico} (IRE)---covering January~1,~2025
through March~18,~2026 (443 days). Each day, 28,364 articles from 16
Peruvian media outlets are classified by Claude Haiku (Anthropic) on a
continuous 0--100 scale for both political and economic uncertainty.  The
indices combine a normalized intensity component---today's aggregate score
relative to a 90-day rolling baseline---with a breadth component---the
fraction of articles exceeding a threshold---using the formula
$\text{IRP}_t = \text{Intensity}_t \times \text{Breadth}_t^{\,\beta}$.
The indices correctly identify known crisis episodes including the
transporters' strike (October 2025, IRP~=~248), the censure of Cabinet
Minister Jeri (January 2026, IRP~=~231), and the Camisea gas crisis
(March 2026, IRE~=~418).  Daily IRP has mean 72.3 and standard
deviation 42.1; IRE has mean 57.8 and standard deviation 65.3,
reflecting concentrated economic coverage.  Monthly IRE correlates
with the PEN/USD exchange rate ($r = 0.398$), while daily and weekly
correlations are near zero, consistent with forward-looking FX markets.
Data are publicly available via the Qhawarina.pe platform.
\end{abstract}

\noindent\textbf{JEL codes:} C43, C82, E32, H56

\noindent\textbf{Keywords:} political uncertainty, economic uncertainty, news-based
index, LLM classification, Peru, exchange rate

\newpage
\setcounter{page}{1}

% ══════════════════════════════════════════════════════════════════════════════
\section{Introduction}
\label{sec:intro}
% ══════════════════════════════════════════════════════════════════════════════

Peru presents a distinctive challenge for researchers studying political
economy: the country has experienced extraordinary political instability
while maintaining broadly stable macroeconomic fundamentals.  Between 2016
and 2026, Peru had eight presidents, multiple congressional dissolutions,
three presidential impeachment attempts, and the arrest of three
ex-presidents---yet GDP growth remained positive in most years and sovereign
spreads stayed comparatively low.  Understanding when and why political
turmoil translates into economic disruption requires a precise, real-time
measurement of both types of uncertainty separately.

No existing measure satisfies these requirements.  The Caldara and
Iacoviello~(\citeyear{caldara2022}) Geopolitical Risk (GPR) index is global
in scope and misses Peru's highly domestic political conflicts.  The Baker,
Bloom, and Davis~(\citeyear{baker2016}) Economic Policy Uncertainty (EPU)
index does not cover Peru.  The International Country Risk Guide (ICRG) is
monthly, subjectively compiled, and published with a lag.  The Banco Central
de Reserva del Per\'u (BCRP) tracks country risk through the JP~Morgan
EMBIG~Peru spread and CDS spreads---financial market measures that are
forward-looking and reflect global risk appetite as much as domestic
conditions.

This paper presents the \emph{\'Indice de Riesgo Pol\'itico} (IRP) and
\emph{\'Indice de Riesgo Econ\'omico} (IRE) for Peru: two daily,
automated, Peru-specific indices constructed from the near-universe of
Peruvian media output.  The indices are built using a four-step pipeline:
(1) RSS feeds from 16 Peruvian outlets are fetched daily; (2) each
article is classified by an AI language model (Claude Haiku, Anthropic)
on two continuous 0--100 scales measuring political and economic
uncertainty; (3) deterministic post-filter rules correct systematic AI
errors; and (4) the indices are computed using an Intensity~$\times$
Breadth$^{\,\beta}$ formula, combining the severity-weighted aggregate
score with the fraction of media attention devoted to risk topics.

The paper makes four contributions. First, we propose and implement a
two-component index formula that separates the \emph{intensity} of risk
coverage from its \emph{breadth} across the media landscape. This
distinction matters: a single newspaper covering a minor cabinet reshuffle
at score~30 is qualitatively different from all~16 outlets covering a
presidential impeachment attempt at score~90. Our formula, inspired by
\cite{caldara2022} but extended to allow continuous AI scoring, captures
both dimensions with a single elasticity parameter~$\beta$.

Second, we extend the AI-GPR methodology of \cite{iacoviello2026} to a
single developing country with a concentrated media market.  We document
the practical challenges---prompt engineering for Spanish-language
political concepts, post-filter rules for false positives---and provide
complete methodological transparency in the appendices.

Third, we provide 443 days of daily data (January~2025--March~2026) that
correctly capture known crisis episodes.  The Paro Nacional de Transportistas
in October~2025 produced the highest IRP spike in the sample (248.2); the
Censura del Gabinete Jeri in January~2026 reached IRP~=~231; and the
Camisea gas disruption in March~2026 produced IRE~=~418, the highest
economic risk reading.

Fourth, we validate the indices against exchange rate movements and discuss
the conditions under which news-based uncertainty measures should and should
not correlate with financial market prices.  We find that daily and weekly
IRP-FX correlations are near zero ($r = -0.011$ and $r = -0.150$,
respectively), while the monthly IRE-FX correlation is $r = 0.398$---a
pattern consistent with forward-looking FX markets that smooth daily
uncertainty but respond to sustained periods of economic disruption.

The paper proceeds as follows. Section~\ref{sec:lit} reviews related
literature. Section~\ref{sec:data} describes the news corpus and AI
classification. Section~\ref{sec:construction} presents the index
construction methodology. Section~\ref{sec:validation} validates the
indices against known events and financial markets.
Section~\ref{sec:properties} describes index properties.
Section~\ref{sec:applications} discusses applications.
Section~\ref{sec:limitations} discusses limitations.
Section~\ref{sec:conclusion} concludes.

% ══════════════════════════════════════════════════════════════════════════════
\section{Related Literature}
\label{sec:lit}
% ══════════════════════════════════════════════════════════════════════════════

\paragraph{News-based uncertainty indices.}
The foundational contribution is \cite{baker2016}, who construct an
Economic Policy Uncertainty (EPU) index from newspaper keyword counts,
forecast disagreement, and tax code expiration data. Their approach has been
replicated for many countries but has not, to our knowledge, been applied to
Peru. \cite{caldara2022} construct the Geopolitical Risk (GPR) index using
keyword frequency across 10 English-language newspapers; their monthly
index covers global geopolitical risk but is not disaggregated to individual
countries. Our approach follows the same news-based logic but replaces
binary keyword matching with continuous AI scoring, and is explicitly
designed for one country's domestic political landscape.

\paragraph{Text analysis in economics.}
\cite{gentzkow2019} survey the use of text as economic data, discussing
dictionary methods, supervised classification, and unsupervised topic
modeling. They emphasize that the choice of representation matters as
much as the choice of method. \cite{manela2017} use front-page newspaper
articles to construct implied volatility measures. \cite{shapiro2022} build
daily news sentiment indices using supervised classification. We follow
these approaches but use a large language model (LLM) rather than a
pre-specified dictionary or human-labeled training set, following the AI-GPR
methodology proposed by \cite{iacoviello2026}.

\paragraph{Firm-level political risk.}
\cite{hassan2019} measure political risk at the firm level using
earnings call transcripts, constructing measures from bigram counts
associated with political topics. Their continuous scoring approach---assigning
intensity to individual documents rather than binary matches---is
conceptually close to ours, though their setting (US firms, English text)
differs substantially from ours (Peruvian political news, Spanish text).

\paragraph{Political uncertainty and financial markets.}
\cite{pastor2013} develop a theoretical model in which political uncertainty
commands a risk premium. Empirically, \cite{julio2012} find that investment
falls sharply in election years, even controlling for aggregate conditions.
\cite{bloom2009} and \cite{ludvigson2021} study uncertainty shocks more
broadly. Our contribution is to provide a daily empirical proxy for
political uncertainty in Peru that can serve as an input to such analyses.

\paragraph{Peru-specific context.}
The BCRP publishes monthly analyses of country risk through the EMBIG spread,
but no daily news-based political or economic uncertainty measure exists for
Peru.  The volatility of Peruvian politics---seven presidents in seven
years, the self-coup attempt of Pedro Castillo in December~2022, continued
confrontation between executive and legislative branches---makes Peru an
especially interesting case study for the political-economic risk literature.

% ══════════════════════════════════════════════════════════════════════════════
\section{Data}
\label{sec:data}
% ══════════════════════════════════════════════════════════════════════════════

\subsection{News Sources}
\label{subsec:sources}

The index is built from 16 Peruvian media outlets monitored via RSS feeds
and historical archives.  The outlets span a range of editorial orientations
and geographic coverage:

\begin{itemize}
  \item \emph{National generalist}: El Comercio, La Rep\'ublica, RPP,
        Correo, Per\'u21
  \item \emph{Business and financial}: Gesti\'on
  \item \emph{Political and analytical}: Caretas, Canal~N
  \item \emph{Official government}: Andina
  \item \emph{Regional}: El B\'uho, Inforegi\'on
  \item \emph{Popular}: Trome, Diario Uno
  \item \emph{Broadcast}: ATV, Panamericana
  \item \emph{Other}: La Raz\'on
\end{itemize}

The sample covers January~1,~2025 through March~18,~2026 (443 days), with
28,364 total articles. Coverage is highest from mid-2025 onward as the
automated ingestion pipeline reached its full set of sources.
Table~\ref{tab:sources} presents source-level statistics.

\begin{table}[h!]
\centering
\caption{Summary Statistics by Source}
\label{tab:sources}
\begin{tabular}{lrrrrrr}
\toprule
Source & N articles & Days active & Mean daily & Mean pol score$^a$ & Mean eco score$^a$ & \% pol$>$20 \\
\midrule
Gesti\'on     & 8,977 & 409 & 21.9 & 26.1 & 21.4 & 28.6 \\
Correo        & 8,389 & 444 & 18.9 & 31.0 & 22.6 & 55.8 \\
El Comercio   & 5,021 & 319 & 15.7 & 28.4 & 22.9 & 26.6 \\
La Rep\'ublica & 2,523 & 126 & 20.0 & 29.4 & 25.3 & 34.6 \\
RPP           & 1,159 & 125 &  9.3 & 27.9 & 18.9 & 13.5 \\
Andina        &   663 &  13 & 51.0 & 13.5 & 22.4 &  6.0 \\
Per\'u21      &   502 &  76 &  6.6 & 26.2 & 23.2 & 25.5 \\
Caretas       &   255 &   7 & 36.4 & 22.2 & 26.6 & 24.3 \\
Canal~N       &   189 &  43 &  4.4 & 27.2 & 25.2 & 37.6 \\
El B\'uho     &   176 &  38 &  4.6 & 26.8 & 24.6 & 41.5 \\
La Raz\'on    &   129 &  56 &  2.3 & 30.0 & 15.3 & 39.5 \\
Panamericana  &   105 &  48 &  2.2 & 27.9 & 22.4 & 40.0 \\
ATV           &    99 &  16 &  6.2 & 15.1 & 16.6 &  6.1 \\
Diario Uno    &    80 &   8 & 10.0 & 17.4 & 17.5 & 16.2 \\
Inforegi\'on  &    76 &  28 &  2.7 & 28.5 & 18.2 & 22.4 \\
Trome         &    21 &  16 &  1.3 & 28.9 & 25.0 & 33.3 \\
\bottomrule
\multicolumn{7}{l}{\footnotesize $^a$ Mean score among articles with score $>0$.}
\end{tabular}
\end{table}

Three sources---Gesti\'on (31.6\% of articles), Correo (29.6\%), and El
Comercio (17.7\%)---together account for 79\% of the corpus.  This
concentration is a structural feature of the Peruvian media market
(Gesti\'on is the only dedicated economic daily) rather than a sampling
choice.  Section~\ref{subsec:limitations_sources} discusses the implications
for index interpretation.

\subsection{AI Classification}
\label{subsec:classification}

Each article is classified by Claude Haiku (model version
\texttt{claude-haiku-4-5}) using two independent calls---one for
political risk and one for economic risk---each with a dedicated system
prompt designed to elicit a 0--100 continuous severity score.  The full
prompts are reproduced in Appendix~\ref{app:prompts}.

The political risk prompt instructs the model to assess whether the
article ``signals that Peru's political institutional trajectory is harder
to predict.'' The economic risk prompt asks whether the article indicates
``active disruption or threat to Peru's economic activity.''  Both prompts
include explicit examples of zero-score articles (sports, weather,
entertainment, irrelevant foreign news) and calibrated examples of
high-score articles (presidential vacancy motions at 80--90; Camisea gas
crisis at 80--90).  The model responds with a single integer score.

The scoring rubric is as follows:
\begin{itemize}
  \item 0: No political (or economic) risk signal
  \item 1--20: Weak signal (routine campaign news, minor bureaucratic friction)
  \item 21--40: Moderate tension (ongoing investigations, sectoral disputes)
  \item 41--60: Significant instability (interpellations, power confrontations)
  \item 61--80: Major crisis (vacancia, censures, national protests)
  \item 81--100: Emergency (consummated vacancy, coup, presidential arrest)
\end{itemize}

\paragraph{Post-classification filtering.}
The AI classifier makes systematic errors that require deterministic
correction.  A \texttt{post\_filter\_scores()} function applies rule-based
fixes after each classification run.  The rules address five failure modes:

\begin{enumerate}
  \item \emph{Sports false positives}: Sports articles (Liga~1, Copa
    Libertadores, NBA, Formula~1, Olympics) receive non-zero political
    scores. The post-filter applies a regex covering all
    Peruvian Liga~1 clubs, European leagues, and international competitions,
    zeroing both scores.

  \item \emph{Entertainment/celebrity}: Tabloid content, music concerts,
    awards ceremonies, and reality TV shows receive false economic signals.
    These are zeroed unconditionally.

  \item \emph{Foreign domestic politics}: Coverage of Argentine, Colombian,
    or US domestic political events that lack direct Peru relevance is
    zeroed. A crisis guard prevents over-aggressive zeroing when a foreign
    event has explicit Peru economic consequences (e.g., US tariffs on
    Peruvian exports).

  \item \emph{Political inversion}: Articles about judicial proceedings
    against political figures are systematically assigned high economic
    scores (eco~=~60--70) and zero political scores by the classifier. The
    post-filter zeros the economic score and sets a minimum political floor
    of 65 for articles matching judicial/political confrontation patterns.

  \item \emph{Score calibration}: Key events known to be systematically
    under-scored are given minimum floor values: presidential vacancy
    motions (pol~$\geq$~90), premier removal (pol~$\geq$~65), minister
    censure (pol~$\geq$~55), Camisea gas crisis (eco~$\geq$~80), and
    Petroper\'u bailout events (eco~$\geq$~75).
\end{enumerate}

The complete list of rules with all regex patterns is in
Appendix~\ref{app:postfilter}.

\paragraph{Classification outcomes.}
Of the 28,364 articles in the corpus, 45.0\% receive a political score of
exactly zero; 84.9\% receive an economic score of exactly zero.  This
asymmetry reflects two facts: Peruvian media covers political news more
broadly than economic instability news, and the economic classifier is
intentionally strict---routine growth statistics, positive economic
announcements, and corporate earnings are scored zero.  Among articles
with non-zero scores, 30.6\% of articles have political score~$>$~20 and
5.9\% have economic score~$>$~20.  Table~\ref{tab:score_dist} shows the
full score distribution.

\begin{table}[h!]
\centering
\caption{Distribution of AI-Assigned Scores}
\label{tab:score_dist}
\small
\begin{tabular}{lrrrr}
\toprule
Score bin & N (political) & \% (political) & N (economic) & \% (economic) \\
\midrule
0         & 12,763 & 45.0 & 24,067 & 84.9 \\
1--10     &  3,685 & 13.0 &  1,176 &  4.1 \\
11--20    &  3,229 & 11.4 &  1,451 &  5.1 \\
21--30    &  2,682 &  9.5 &    778 &  2.7 \\
31--40    &  2,362 &  8.3 &    306 &  1.1 \\
41--50    &  1,472 &  5.2 &    191 &  0.7 \\
51--60    &    793 &  2.8 &    212 &  0.7 \\
61--70    &  1,125 &  4.0 &     99 &  0.3 \\
71--80    &    121 &  0.4 &     84 &  0.3 \\
81--90    &    131 &  0.5 &      0 &  0.0 \\
91--100   &      1 &  0.0 &      0 &  0.0 \\
\midrule
Total     & 28,364 & 100.0 & 28,364 & 100.0 \\
\bottomrule
\end{tabular}
\end{table}

The score distributions show a clear right-skew with mass at zero: most
articles are irrelevant to political or economic risk, and genuinely
high-risk articles (score~$>$~60) are rare.  The economic distribution is
more concentrated at zero, reflecting the narrower scope of economic
instability news relative to general political coverage.

\subsection{Comparison with Keyword Approaches}
\label{subsec:keyword_comparison}

The Caldara and Iacoviello (\citeyear{caldara2022}) GPR index uses binary
keyword matching: an article either matches one of their risk-associated
keywords or it does not.  Our approach differs in three ways.

First, AI scoring provides context-sensitivity that keyword matching
cannot.  The word ``crisis'' appears in sports articles (``crisis of
confidence in the attacking line''), economic opinion columns (``the
fiscal framework is in crisis''), and in descriptions of genuine political
emergencies.  A keyword counter cannot distinguish these; a language model,
given appropriate prompting, can.  Similarly, Peru-specific institutional
concepts---\emph{voto de confianza}, \emph{vacancia presidencial},
\emph{paro nacional}---require prompts calibrated to the Peruvian context,
not a generic Spanish political vocabulary.

Second, continuous scoring (0--100) captures intensity, not just
occurrence.  A day when 40 articles each score 30 is qualitatively
different from a day when 4 articles each score 90.  The keyword approach
treats both as equivalent counts.  Our formula uses the sum of scores
(raw intensity) as the numerator, so higher-severity articles contribute
proportionally more.

Third, and most honestly, AI scoring introduces new failure modes absent
from keyword matching.  The classifier makes systematic errors---sports
false positives, political inversion---that require post-hoc rules to
correct.  Keyword matching has more predictable failure modes.  We do not
claim that AI scoring is universally superior; we claim it is better suited
to the specific problem of measuring uncertainty in a single country's
Spanish-language media.

% ══════════════════════════════════════════════════════════════════════════════
\section{Index Construction}
\label{sec:construction}
% ══════════════════════════════════════════════════════════════════════════════

\subsection{Two-Component Formula}
\label{subsec:formula}

The index construction follows a four-step procedure for each day~$t$.

\paragraph{Step 1: Raw intensity.}
Let $s_{it}^{\text{pol}} \in [0, 100]$ denote the political score of article $i$
on day $t$, and let $N_t$ denote the total number of articles on day $t$.  The
raw political intensity is:
\begin{equation}
  \text{Raw}_t^{\text{pol}} = \sum_{i=1}^{N_t} s_{it}^{\text{pol}}
  \label{eq:raw}
\end{equation}
and analogously for $\text{Raw}_t^{\text{eco}}$.  This is the sum---not the
mean---of all article scores.  A day with 12 outlets each covering a major
crisis at score~80 produces a raw intensity of $12 \times 80 = 960$,
correctly reflecting the uniformity of media attention.  This mirrors the
GPR approach, where the aggregate count of matching articles is the
numerator.

\paragraph{Step 2: Normalized intensity.}
Raw intensity is normalized by a rolling baseline to produce a
mean-100 index:
\begin{equation}
  I_t^{\text{pol}} = \frac{\text{Raw}_t^{\text{pol}}}{\bar{B}_t^{\text{pol}}} \times 100
  \label{eq:intensity}
\end{equation}
where the baseline is a rolling mean over the past $W$ days:
\begin{equation}
  \bar{B}_t^{\text{pol}} = \frac{1}{\min(W, t)} \sum_{j=1}^{\min(W, t)} \text{Raw}_{t-j}^{\text{pol}}
  \label{eq:baseline}
\end{equation}
with $W = 90$ (the default baseline window).  During the first 90 days of
the sample, the baseline is an expanding mean over all available history
(\texttt{min\_periods=1}).  After 90 days, the baseline is a fixed 90-day
rolling mean.  The baseline normalization means $I_t^{\text{pol}} = 100$
on a typical day, $I_t^{\text{pol}} = 200$ on a day twice as intense as
the average, and so on.

\paragraph{Step 3: Breadth.}
The breadth component measures the fraction of all media articles that
are devoted to political uncertainty:
\begin{equation}
  D_t^{\text{pol}} = \frac{\#\{i : s_{it}^{\text{pol}} \geq \tau\}}{N_t}
  \label{eq:breadth}
\end{equation}
where $\tau = 20$ is the threshold score for an article to ``count'' as a
political risk article.  Articles below this threshold are treated as
background noise: routine political coverage that does not signal
instability.  This component is directly analogous to the GPR breadth
measure---the share of media attention devoted to the topic of interest.

\paragraph{Step 4: Combined index.}
The IRP and IRE are:
\begin{equation}
  \text{IRP}_t = I_t^{\text{pol}} \times \left(D_t^{\text{pol}}\right)^{\beta_{\text{pol}}}
  \label{eq:irp}
\end{equation}
\begin{equation}
  \text{IRE}_t = I_t^{\text{eco}} \times \left(D_t^{\text{eco}}\right)^{\beta_{\text{eco}}}
  \label{eq:ire}
\end{equation}
with $\beta_{\text{pol}} = 0.5$ and $\beta_{\text{eco}} = 0.3$.

The exponent $\beta$ governs how strongly breadth amplifies or dampens the
intensity signal.  With $\beta = 0.5$: breadth~=~0.25 gives multiplier
$0.25^{0.5} = 0.50$; breadth~=~0.50 gives multiplier $0.50^{0.5} = 0.71$.
The exponent is set below 1 so that breadth scales the index without
dominating it: a genuine high-intensity day with moderate breadth should
still produce a spike.

The rationale for $\beta_{\text{eco}} < \beta_{\text{pol}}$ is structural:
economic news in Peru is concentrated in specialized outlets (Gesti\'on,
the economics section of El Comercio).  Even on major economic crisis days,
$D_t^{\text{eco}}$ rarely exceeds 10--15\% of all articles, while
$D_t^{\text{pol}}$ can reach 30--40\% during political crises.  A lower
$\beta_{\text{eco}}$ prevents the breadth multiplier from systematically
suppressing the economic index.  Empirically, the mean political breadth in
the sample is 20.3\% while the mean economic breadth is 3.5\%.

\subsection{Temporal Smoothing}
\label{subsec:smoothing}

A crisis does not disappear from the information environment the day after
a major event: follow-up reporting, analysis, and political responses
continue for several days.  To capture this persistence, we apply
exponential moving average (EMA) smoothing:
\begin{equation}
  \text{IRP}_t^{\text{smooth}} = \alpha \cdot \text{IRP}_{t-1}^{\text{smooth}}
    + (1 - \alpha) \cdot \text{IRP}_t
  \label{eq:ema}
\end{equation}
with $\alpha = 0.3$, giving 70\% weight on the current day and 30\%
persistence from recent history.  In the pandas implementation, this
corresponds to \texttt{ewm(alpha=0.7, adjust=False)}.  The smoothed
series is reported alongside the raw series for all analyses.

\subsection{Parameter Sensitivity}
\label{subsec:sensitivity}

Table~\ref{tab:sensitivity} reports the average IRP and IRE for the week of
March~14--18,~2026 under alternative parameter choices.  This week is
informative because it contains both a political event (voto de confianza,
March~18) and ongoing economic disruption (Camisea gas crisis), so both
indices are above their sample averages.

\begin{table}[h!]
\centering
\caption{Parameter Sensitivity Analysis: Average IRP and IRE, March~14--18, 2026}
\label{tab:sensitivity}
\small
\begin{tabular}{ccccc}
\toprule
$\beta_{\text{pol}}$ & $\beta_{\text{eco}}$ & Window $W$ (days) & Avg.\ IRP & Avg.\ IRE \\
\midrule
\multicolumn{5}{l}{\textit{Varying $\beta_{\text{pol}}$ (baseline: $\beta_{\text{eco}}=0.3$, $W=90$)}} \\
0.3 & 0.3 & 90 & 183.8 & 120.2 \\
\textbf{0.5} & \textbf{0.3} & \textbf{90} & \textbf{129.2} & \textbf{120.2} \\
0.7 & 0.3 & 90 &  91.0 & 120.2 \\
\midrule
\multicolumn{5}{l}{\textit{Varying $\beta_{\text{eco}}$ (baseline: $\beta_{\text{pol}}=0.5$, $W=90$)}} \\
0.5 & 0.2 & 90 & 129.2 & 163.2 \\
0.5 & \textbf{0.3} & 90 & 129.2 & \textbf{120.2} \\
0.5 & 0.5 & 90 & 129.2 &  65.3 \\
\midrule
\multicolumn{5}{l}{\textit{Varying baseline window $W$ (baseline: $\beta_{\text{pol}}=0.5$, $\beta_{\text{eco}}=0.3$)}} \\
0.5 & 0.3 &  60 & 104.0 &  94.9 \\
0.5 & 0.3 & \textbf{90}  & \textbf{129.2} & \textbf{120.2} \\
0.5 & 0.3 & 120 & 138.5 & 143.4 \\
\bottomrule
\multicolumn{5}{l}{\footnotesize Bold row is the baseline specification.}\\
\multicolumn{5}{l}{\footnotesize Daily detail (baseline): Mar-14: IRP=61.1 IRE=121.4; Mar-15: IRP=37.5 IRE=30.2;}\\
\multicolumn{5}{l}{\footnotesize Mar-16: IRP=123.8 IRE=169.8; Mar-17: IRP=197.2 IRE=116.7; Mar-18: IRP=226.2 IRE=163.0.}
\end{tabular}
\end{table}

Three observations from Table~\ref{tab:sensitivity}.  First, IRP is
monotonically decreasing in $\beta_{\text{pol}}$: a higher exponent
amplifies the breadth penalty, compressing the index toward zero on days
with moderate breadth.  Second, IRE is highly sensitive to $\beta_{\text{eco}}$:
because economic breadth is structurally low, increasing $\beta_{\text{eco}}$
from 0.3 to 0.5 cuts the average IRE nearly in half.  This confirms the
choice of a lower $\beta$ for the economic dimension.  Third, a longer
baseline window $W$ raises both indices on this particular week---because
the baseline is computed over a calmer period, so the same raw intensity
appears larger relative to history.  The qualitative pattern (IRP rising
toward the March~18 voto de confianza, IRE elevated throughout by the energy
crisis) is robust to all parameter choices.

\subsection{Relationship to Caldara and Iacoviello (2022)}
\label{subsec:ci_comparison}

Table~\ref{tab:comparison} summarizes the methodological relationship
between the GPR and the IRP/IRE.

\begin{table}[h!]
\centering
\caption{Comparison of GPR (Caldara \& Iacoviello, 2022) and IRP/IRE (this paper)}
\label{tab:comparison}
\small
\begin{tabularx}{\textwidth}{lXX}
\toprule
Feature & GPR (C\&I 2022) & IRP/IRE (this paper) \\
\midrule
Scope & Global geopolitical & Peru domestic \\
Frequency & Monthly (daily available) & Daily \\
Classification method & Binary keyword match & Continuous AI score (0--100) \\
Formula & Share of matching articles & Intensity $\times$ Breadth$^{\,\beta}$ \\
Risk dimensions & Single (geopolitical) & Two (political + economic) \\
Sources & 10 English newspapers & 16 Peruvian Spanish outlets \\
Normalization & Mean = 100 (1985--2019 baseline) & Mean = 100 (90-day rolling baseline) \\
Backward coverage & From 1899 (newspaper archives) & January 2025 onward \\
\bottomrule
\end{tabularx}
\end{table}

Our approach differs from \cite{caldara2022} in three important ways.
We use a rolling 90-day baseline rather than a fixed historical baseline,
so the index reflects \emph{current} conditions relative to recent history
rather than relative to a 30-year normal.  We use separate dimensions for
political and economic risk, capturing the common finding that political
and economic disruptions in Peru often have different causes and different
transmission mechanisms.  And we use AI classification, which allows
context-sensitive scoring of Spanish-language political concepts that have
no direct English keyword equivalent.

% ══════════════════════════════════════════════════════════════════════════════
\section{Validation}
\label{sec:validation}
% ══════════════════════════════════════════════════════════════════════════════

\subsection{Event Detection}
\label{subsec:events}

The most direct validation is to check whether the indices spike around
known political and economic events.  Table~\ref{tab:events} reports, for
each major event in the sample period, the date with the highest IRP within
a $\pm$3-day window, together with the index values and coverage statistics
on that date.

\begin{table}[h!]
\centering
\caption{Known Events and Index Response}
\label{tab:events}
\small
\begin{tabular}{lrrrrrrr}
\toprule
Event & Date & IRP & IRP$_s$ & IRE & IRE$_s$ & Articles & $D^{\text{pol}}$ \\
\midrule
Cerr\'on Desaparece (pr\'ofugo) & Mar-18, 2025 & 127.6 & 115.5 & 16.9 & 26.8 & 32 & 0.719 \\
Impunidad Boluarte & Apr-23, 2025 & 82.6 & 66.7 & 9.1 & 11.5 & 28 & 0.714 \\
Vizcarra Prisi\'on Preventiva & Jun-16, 2025 & 137.8 & 119.4 & 0.0 & 5.1 & 47 & 0.532 \\
Fragmentaci\'on Izquierda & Jul-17, 2025 & 88.8 & 78.2 & 96.1 & 88.6 & 39 & 0.615 \\
Colapso Ministerial Boluarte & Aug-19, 2025 & 145.5 & 127.0 & 112.6 & 80.5 & 57 & 0.596 \\
Paro Transportistas & Oct-10, 2025 & 248.2 & 225.0 & 150.2 & 128.0 & 74 & 0.730 \\
Vacancia Boluarte & Oct-22, 2025 & 153.2 & 131.5 & 73.6 & 84.2 & 83 & 0.542 \\
Castillo Condenado & Nov-13, 2025 & 112.8 & 89.4 & 147.6 & 115.2 & 66 & 0.530 \\
Censura Gabinete Jeri & Jan-21, 2026 & 231.1 & 201.9 & 66.3 & 66.3 & 127 & 0.504 \\
Paro Transportistas II & Feb-04, 2026 & 91.5 & 84.9 & 37.1 & 43.8 & 111 & 0.351 \\
Crisis Energ\'etica Camisea & Mar-12, 2026 & 138.9 & 118.9 & 315.3 & 292.7 & 551 & 0.176 \\
Voto Confianza Miralles & Mar-18, 2026 & 226.2 & 208.9 & 163.0 & 151.4 & 932 & 0.186 \\
\bottomrule
\multicolumn{8}{l}{\footnotesize IRP$_s$, IRE$_s$ = EMA-smoothed values. $D^{\text{pol}}$ = political breadth.}
\end{tabular}
\end{table}

The indices correctly identify all 12 events listed.  Several patterns
merit discussion.  First, the Paro Nacional de Transportistas in October~2025
produced the highest IRP reading in the sample (IRP~=~248.2), with 73\% of
articles on that day scoring above the threshold---the highest breadth
recorded.  This was a genuine national crisis: road blockades nationwide,
supply chain disruption, and a government that nearly fell.  Second, the
Censura Jeri episode (January~2026) reached IRP~=~231.1 with 50.4\% breadth,
reflecting broad coverage of a ministerial censure that threatened the
cabinet's survival.  Third, the Camisea gas crisis in March~2026 produced
the highest single-day IRE reading (IRE~=~418.2 on March~3, with the
event table showing 315.3 on March~12), reflecting the catastrophic economic
impact of 13 consecutive days without natural gas to industry and households.
The March~12 Camisea reading has lower political breadth (0.176)
than the political events, confirming that the two indices measure distinct
phenomena.

Fourth, the voto de confianza for Cabinet Miralles (March~18,~2026) appears
simultaneously in both indices: IRP~=~226.2 and IRE~=~163.0, with 932 total
articles on that single day---the largest single-day article volume in the
sample.  This confirms that political uncertainty can spillover into
economic uncertainty when the political event involves fiscal or energy policy.

Fifth, note the clear IRP/IRE divergence for the Vizcarra Prisi\'on
Preventiva event (June~2025): IRP~=~137.8 but IRE~=~0.0, correctly
identifying this as a purely political event (the imprisonment of a former
president) with no contemporaneous economic disruption.

\subsection{Correlation with Financial Markets}
\label{subsec:fx}

If the IRP/IRE measure genuine uncertainty, they should correlate positively
with measures of country risk at sufficiently low frequencies.  We use the
PEN/USD exchange rate as the primary financial benchmark.  A higher
PEN/USD means more soles per dollar (peso depreciation), which typically
reflects increased country risk perception.

Table~\ref{tab:fx} reports Pearson correlations of IRP and IRE with the
daily change in PEN/USD, at daily, weekly, and monthly aggregation.

\begin{table}[h!]
\centering
\caption{Correlation of IRP/IRE with PEN/USD Daily Change}
\label{tab:fx}
\begin{tabular}{lrrrrrr}
\toprule
Frequency & $r(\text{IRP})$ & $p$-value & $r(\text{IRE})$ & $p$-value & $N$ \\
\midrule
Daily   & $-0.011$ & 0.852 & $0.113$ & 0.051 & 299 \\
Weekly  & $-0.150$ & 0.236 & $0.217$ & 0.084 &  64 \\
Monthly & $-0.047$ & 0.867 & $0.398$ & 0.142 &  15 \\
\bottomrule
\multicolumn{6}{l}{\footnotesize FX data from Banco Central de Reserva del Per\'u. 299 overlapping trading days.} \\
\multicolumn{6}{l}{\footnotesize Monthly correlation uses 15 month-average observations.}
\end{tabular}
\end{table}

Three findings emerge.  First, daily IRP has essentially zero correlation
with daily FX changes ($r = -0.011$, $p = 0.852$).  This is expected and
does not invalidate the index.  The exchange rate reflects all available
information about future risk, not just today's news; BCRP intervention
further smooths daily moves; and measurement noise in the daily index is
substantial when $N_t$ is small (median 47 articles per day in early
months).

Second, the monthly IRE correlation with monthly FX changes is $r = 0.398$
($p = 0.142$; 15 observations).  Though not statistically significant at
conventional thresholds given only 15 months of data, the magnitude is
economically meaningful and directionally consistent with the hypothesis
that sustained economic disruption puts sustained pressure on the sol.

Third, the IRE-FX correlation is uniformly higher than the IRP-FX
correlation across all frequencies.  This is plausible: the exchange rate
is more directly linked to economic disruption (supply chain interruptions,
fiscal stress, investment flight) than to political events per se.  Peru's
BCRP has a long track record of intervening to smooth politically-driven
FX volatility, which would attenuate the IRP-FX link.

\subsection{Comparison with EMBIG Peru}
\label{subsec:embig}

EMBIG Peru---the JP~Morgan Emerging Market Bond Index spread for Peru---is a
daily market-based measure of sovereign risk that incorporates both political
and economic information.  We do not have daily EMBIG data for the full
sample period.  The BCRP publishes monthly averages.  Based on those data,
EMBIG Peru averaged approximately 120--130 basis points throughout 2025,
with a modest widening in October--November (coinciding with the
transporters' strike and Castillo sentencing) and again in March~2026
(energy crisis).  This broad co-movement is consistent with our indices.
A formal correlation analysis awaits a longer time series.

\subsection{False Positive Audit}
\label{subsec:fp_audit}

A formal false positive audit (randomly sampling 100 articles per score
quintile and comparing AI scores to human judgment) has not yet been
conducted.  Based on qualitative review of index spikes and contributing
articles, the most common false positive types are:

\begin{enumerate}
  \item \emph{Electoral coverage misclassified as crisis}: Articles about
    campaign announcements, opinion polls, or candidate debates
    occasionally receive political scores of 50--70 when the appropriate
    score is 15--25.  The post-filter applies a 15--25 floor for campaign
    news, partially addressing this.

  \item \emph{Foreign economic news affecting Peru-adjacent sectors}:
    Articles about US trade policy or commodity price movements sometimes
    receive high economic scores when the Peru-specific impact is unclear.
    The post-filter crisis guard preserves these when explicit Peru impact
    is mentioned.

  \item \emph{Opinion column inflation}: Economic analyst opinion columns
    receive scores of 40--60 when the appropriate score for ``one economist's
    view'' is 15--20.  The post-filter caps named-economist opinion columns
    at 20.
\end{enumerate}

A formal audit is planned for future work; see Section~\ref{sec:limitations}.

% ══════════════════════════════════════════════════════════════════════════════
\section{Properties of the Index}
\label{sec:properties}
% ══════════════════════════════════════════════════════════════════════════════

\subsection{Descriptive Statistics}
\label{subsec:descriptive}

Table~\ref{tab:summary} reports summary statistics for daily IRP and IRE
over the full 443-day sample.

\begin{table}[h!]
\centering
\caption{Summary Statistics of Daily IRP and IRE (January 2025--March 2026)}
\label{tab:summary}
\begin{tabular}{lrrrrrrrrrr}
\toprule
Series & Mean & Median & Std & Min & Max & Skew & Kurt & $\rho_1$ & $\rho_7$ & $\rho_{30}$ \\
\midrule
IRP (raw)      & 72.3 & 65.4 & 42.1 &  7.8 & 302.6 & 1.63 & 4.72 & 0.456 & 0.240 & $-0.060$ \\
IRP (smoothed) & 72.3 & 67.6 & 35.4 & 18.4 & 287.8 & 1.76 & 5.81 & 0.656 & 0.249 & $-0.021$ \\
IRE (raw)      & 57.8 & 38.7 & 65.3 &  0.0 & 418.2 & 2.35 & 7.17 & 0.510 & 0.347 & $-0.015$ \\
IRE (smoothed) & 57.8 & 41.4 & 56.6 &  0.0 & 355.3 & 2.41 & 7.40 & 0.730 & 0.415 & $-0.018$ \\
\bottomrule
\multicolumn{11}{l}{\footnotesize $\rho_k$ = sample autocorrelation at lag $k$ days. $N = 443$ days.}
\end{tabular}
\end{table}

Both series are right-skewed with excess kurtosis, consistent with rare but
large crisis episodes.  The IRP has positive skewness (1.63) and kurtosis
(4.72), indicating a fat right tail.  The IRE is more extreme: skewness 2.35
and kurtosis 7.17.  The IRE all-time maximum of 418.2 (on March~3,~2026,
during the Camisea gas crisis) represents a day with raw economic intensity
more than four times the 90-day average, with breadth_eco~=~0.060.  The
large kurtosis reflects these tail events.

The lag-1 autocorrelation ($\rho_1$) is moderate for IRP (0.456) and higher
for IRE (0.510), suggesting that risk episodes persist for a few days.
The lag-7 autocorrelation drops substantially ($\rho_7^{\text{IRP}} = 0.240$,
$\rho_7^{\text{IRE}} = 0.347$), and lag-30 autocorrelation is near zero.
This pattern is consistent with crises that last days to weeks rather than
months.  The smoothed series has higher $\rho_1$ by construction (EMA
introduces mechanical persistence) and similar $\rho_7$.

The IRP mean of 72.3 is below 100.  This reflects the rolling
baseline normalization during the early sample months: the baseline window
was short (expanding mean over fewer than 90 days), so as the sample grew
and more high-risk days were added to the rolling history, the baseline
rose, compressing recent index values.  As the sample extends, the mean
will converge toward 100 by construction.

\subsection{Persistence of High-Risk Episodes}
\label{subsec:persistence}

We define a high-risk episode as a period of consecutive days with
IRP~$\geq$~150.  Over the 443-day sample, there are 21 days with
IRP~$\geq$~150, spanning 14 distinct episodes with average duration 1.5
days.  The longest episode was 4 consecutive days.  This is brief: political
crises in Peru resolve quickly (through cabinet reshuffles, Congressional
votes, or simply news fatigue) or escalate further.

For IRE, there are 33 days above 150 in 22 distinct episodes, also
averaging 1.5 days.  The energy crisis of February--March 2026 produced the
most sustained elevated IRE: multiple days above 200 during the Camisea
disruption.

\subsection{Seasonality and Periodicity}
\label{subsec:seasonality}

IRP shows a clear day-of-week pattern.  Wednesday has the highest average
IRP (91.3), followed by Thursday (85.9) and Monday (78.8); Saturday has the
lowest (36.3), followed by Sunday (60.4).  This reflects two facts: fewer
articles are published on weekends (the breadth component falls), and
political events (Congressional sessions, cabinet votes, press conferences)
concentrate in weekdays.

IRP does not show a strong intra-year seasonal pattern given our 15-month
sample---too short to identify annual seasonality reliably.  The January~2026
and March~2026 spikes reflect specific political and economic events rather
than calendar effects.

As Peru approaches its April~2026 presidential elections, IRP has trended
upward since January~2026.  Average IRP in the first quarter of 2026
(through March~18) is 85.6, compared to a monthly range of 50--90 during
most of 2025.  Whether this reflects genuine electoral uncertainty or
mechanical effects from increased political news volume requires more data.

\subsection{Political vs.\ Economic Correlation}
\label{subsec:irp_ire_corr}

The daily correlation between IRP and IRE is $r = 0.181$ ($p < 0.001$),
indicating a positive but weak association.  The monthly correlation is
$r = 0.242$ ($p = 0.386$; 15 observations), which is not statistically
significant given the short sample.

The two indices co-move during events that generate both political and
economic uncertainty simultaneously---general strikes (where supply
disruption is economic, but the decision to strike is political),
and the March 2026 energy-political crisis (where the Camisea gas disruption
was economic, but the political management and voto de confianza created
concurrent political uncertainty).

The indices diverge in the more common case where one type of uncertainty
dominates.  The Vizcarra imprisonment (June~2025) is purely political
(IRP~=~138, IRE~=~0).  The Castillo sentencing (November~2025) has moderate
IRP (112.8) but higher IRE (147.6)---because the sentencing reopened
questions about the social programs Castillo had promised, creating economic
uncertainty about redistribution.  The Camisea gas crisis has the most
extreme divergence: IRE~=~418.2 versus IRP~=~138.9 on March~3.

% ══════════════════════════════════════════════════════════════════════════════
\section{Applications}
\label{sec:applications}
% ══════════════════════════════════════════════════════════════════════════════

The IRP and IRE are designed with several use cases in mind:

\paragraph{Central bank monitoring.}
The BCRP monitors political risk through financial market indicators (EMBIG,
CDS spreads, PEN/USD).  These are forward-looking and reflect global risk
appetite as much as domestic conditions.  The IRP/IRE provide a
complementary text-based measure that captures domestic political content
directly.  A sustained IRP~$\geq$~150 episode could serve as an early warning
for the Directorio's risk assessment.

\paragraph{Investment decisions.}
Foreign investors in Peruvian mining, energy, or infrastructure assets need
real-time measures of the policy environment.  IRP captures whether the
Congress-Executive relationship is deteriorating (cabinet censures, voto de
confianza tensions) before it manifests in permitting delays or contract
disputes.  IRE captures sectoral disruptions (mining blockades, gas
supply interruptions) before they show up in production data.

\paragraph{Academic research.}
The IRP and IRE can serve as control variables in event studies, regression
discontinuity designs, or time series models.  For example, one could
estimate the effect of electoral uncertainty (IRP) on formal employment
during the April~2026 campaign period, controlling for concurrent economic
disruption (IRE).  The daily frequency and precise measurement allow
identification strategies that monthly measures preclude.

\paragraph{Journalistic context.}
The indices provide quantitative context for news coverage:
``The voto de confianza on March~18 occurred on a day with IRP~=~226, the
second-highest reading since October~2025.'' This contextualizes individual
events within the broader uncertainty landscape.

% ══════════════════════════════════════════════════════════════════════════════
\section{Limitations and Future Work}
\label{sec:limitations}
% ══════════════════════════════════════════════════════════════════════════════

We discuss six limitations honestly.

\paragraph{1. AI classification quality.}
Claude Haiku makes systematic errors that the post-filter rules address
only partially.  Sports articles, celebrity news, and foreign domestic
politics were the most common false positives identified during qualitative
review.  The post-filter regex patterns are necessarily incomplete: they
cover known failure modes but will miss novel patterns.  A more robust
solution---either a fine-tuned model trained on human-labeled Peruvian
news, or a two-stage filter that screens for relevance before scoring
severity---would improve quality.  We have not conducted a formal false
positive audit to quantify error rates by score quintile.

\label{subsec:limitations_sources}
\paragraph{2. Source concentration.}
Three outlets---Gesti\'on (31.6\%), Correo (29.6\%), and El Comercio
(17.7\%)---account for 79\% of the corpus.  The index is thus most sensitive
to the editorial choices of these three outlets.  If, for example, Correo
systematically amplifies political conflict relative to other outlets, the
IRP will be biased upward.  More concerning is that Gesti\'on's dominance
in economic coverage means the IRE is heavily influenced by a single
outlet's judgment about what constitutes economic risk.  Source-level
credibility weighting (implemented but disabled in the current pipeline)
would allow researchers to test sensitivity to this concentration.

\paragraph{3. Short history.}
443 days of data is insufficient for many econometric applications.  A
reliable autocorrelation function requires several hundred lags; business
cycle analysis requires multiple cycles; electoral analysis requires
multiple electoral cycles.  Backward extension to 2000 (or earlier) using
keyword-based classification of digitized newspaper archives is planned but
requires access to historical archives not yet integrated into the pipeline.

\paragraph{4. No causal identification.}
The IRP and IRE measure \emph{media-reported} uncertainty, not
\emph{objective} political or economic risk.  Media can amplify genuine
uncertainty, underreport suppressed crises, or create uncertainty through
framing.  We do not estimate the causal effect of IRP on any outcome.
With only 443 days and 4 quarters of economic data, credible
identification of IRP's macroeconomic effects is not possible.
This paper is a measurement contribution, not a causal one.

\paragraph{5. Parameter sensitivity.}
The breadth exponents ($\beta_{\text{pol}} = 0.5$, $\beta_{\text{eco}} = 0.3$),
threshold ($\tau = 20$), and baseline window ($W = 90$) are chosen based on
qualitative assessment and robustness to known events.  They are not
estimated from data.  Table~\ref{tab:sensitivity} shows qualitative
robustness, but different parameter choices produce index levels that differ
by a factor of two in some configurations.  Researchers using the indices
for specific applications should assess sensitivity to the parameters relevant
for their application.

\paragraph{6. LLM model drift.}
Claude Haiku's behavior may change across API versions as Anthropic releases
model updates.  Classification consistency over time is not guaranteed if
the underlying model weights change.  We do not have access to version
control over the model used for historical classifications.  If the model's
calibration changes, the index levels pre- and post-change would not be
comparable.  Recording the exact model version with each classification
run, and periodically re-classifying a fixed reference set of articles to
detect drift, are best practices we recommend for future implementation.

\paragraph{Future work.}
Several extensions are planned:
\begin{enumerate}
  \item Backward extension to 2000 using keyword-based proxies applied to
    digitized archives of El Comercio and La Rep\'ublica.
  \item Formal false positive audit: random sample of 100 articles per score
    quintile, human labeling, and precision/recall by quintile.
  \item Sectoral economic sub-indices: separate IRE components for
    mining, energy, fiscal, and trade.
  \item Cross-country comparison: applying the same methodology to
    Colombia, Chile, and Mexico to benchmark Peru's political volatility.
  \item Formal causal analysis: once the index has sufficient history, local
    projections (\citealt{jorda2005}) to estimate dynamic effects of IRP
    shocks on investment and employment.
\end{enumerate}

% ══════════════════════════════════════════════════════════════════════════════
\section{Conclusion}
\label{sec:conclusion}
% ══════════════════════════════════════════════════════════════════════════════

We have constructed and validated two daily indices---the IRP (political
risk) and IRE (economic risk)---for Peru, covering January~2025 through
March~2026.  The indices combine AI-assigned severity scores with a
breadth measure of media attention into a formula calibrated to Peru's
concentrated media landscape.  They correctly identify all major political
and economic events in the sample, including the transporters' strike
(IRP~=~248), the Jeri censure (IRP~=~231), and the Camisea gas crisis
(IRE~=~418).  Monthly IRE correlates with exchange rate movements
($r = 0.398$), while daily correlations are near zero as expected for
forward-looking financial markets.

The principal limitation is the short sample (443 days).  As the pipeline
continues running, the data will become more useful for time series
analysis, event studies, and eventually causal identification.  The data
are publicly available via the Qhawarina.pe platform and the statistics
are fully reproducible from the script \texttt{paper/compute\_paper\_stats.py}
in the repository.

Our broader contribution is methodological: we document the practical steps
required to adapt the AI-GPR methodology (\citealt{iacoviello2026}) to a
single developing country with a concentrated, Spanish-language media
market.  The post-filter rules, the dual-dimension formula, and the
calibration choices are all fully transparent in the appendices.
Researchers applying similar methods to other countries can build on
or depart from these choices as appropriate for their context.

% ══════════════════════════════════════════════════════════════════════════════
% References
% ══════════════════════════════════════════════════════════════════════════════

\clearpage
\bibliographystyle{agsm}
\bibliography{references}

% ══════════════════════════════════════════════════════════════════════════════
% Appendices
% ══════════════════════════════════════════════════════════════════════════════

\clearpage
\appendix
\titleformat{\section}[block]{\large\bfseries}{Appendix \thesection:}{0.5em}{}

\section{Full Classification Prompts}
\label{app:prompts}

This appendix reproduces the complete system prompts sent to Claude Haiku
for political and economic risk classification.  Both prompts are sent as
\texttt{system} messages; the user message lists the articles to classify
using the template in Section~\ref{app:user_template}.

\subsection*{Political Risk Classification Prompt}

The following prompt is sent as the \texttt{system} message for IRP scoring:

\begin{lstlisting}
Eres un analista de riesgo politico peruano. Tu tarea: evaluar si este
articulo de prensa senala INCERTIDUMBRE politica institucional en el Peru.

NO es riesgo politico:
- Deportes (futbol, voley, tenis, retiros de deportistas, resultados
  deportivos) -> 0
- Farandula, celebridades, entretenimiento, TV -> 0
- Clima, alertas meteorologicas, lluvias, vientos, oleajes -> 0
- Tecnologia, gadgets, aplicaciones, dispositivos medicos -> 0
- Noticias internacionales sin impacto directo en la politica peruana -> 0
- Recetas, turismo, estilo de vida, salud general -> 0
- Obituarios de personas no politicas -> 0

SI es riesgo politico (y debe puntuar ALTO):
- Vacancia presidencial, mocion de vacancia -> 75-90
- Voto de confianza (cualquier mencion) -> 65-80
- Interpelacion o censura de ministros -> 55-70
- Cambio de gabinete bajo presion -> 60-75
- Conflicto entre Ejecutivo y Congreso -> 55-70
- Detencion, arresto, orden de captura de presidente o ex-presidente -> 80+
- Corrupcion de funcionarios activos, allanamientos, casos judiciales -> 55-70
- Asesinato de figura politica, funcionario, asesor -> 65-80
- Crimen organizado infiltrado en instituciones politicas -> 70+
- Fallos del Tribunal Constitucional sobre temas politicos -> 55-65
- Protestas masivas que amenazan continuidad del gobierno -> 65-80
- Estado de emergencia por conflicto social -> 60-70

MODERADO (campanas electorales 2026):
- Candidatos presidenciales haciendo propuestas de campana -> 15-25
- Debates, foros de candidatos -> 20-30
- Encuestas electorales -> 15-25
- Candidato con problemas judiciales activos -> 55-70

Escala:
0 = Sin relevancia politica
1-20 = Senal politica debil (campana rutinaria, burocracia menor)
21-40 = Tension baja-moderada (investigaciones, huelgas locales)
41-60 = Inestabilidad significativa (interpelaciones, conflictos entre poderes)
61-80 = Crisis mayor (vacancia, censuras, protestas nacionales)
81-100 = Emergencia (vacancia consumada, golpe, renuncia presidencial)

REGLA: Si un articulo usa una crisis economica como CONTEXTO para una
maniobra politica, el puntaje POLITICO debe ser alto.
\end{lstlisting}

\subsection*{Economic Risk Classification Prompt}

The following prompt is sent as the \texttt{system} message for IRE scoring:

\begin{lstlisting}
Eres un analista de riesgo economico peruano. Tu tarea: evaluar si este
articulo senala INCERTIDUMBRE sobre la actividad economica del Peru.

NO es riesgo economico:
- Deportes, farandula, entretenimiento -> 0
- Clima rutinario -> 0
- Tecnologia, gadgets, apps -> 0
- Desastres internacionales sin conexion a Peru -> 0
- Datos economicos POSITIVOS o RUTINARIOS (crecimiento normal, inversion
  anunciada, records de exportacion) -> 0
- Noticias de candidatos sobre temas NO economicos -> 0
- Crimen comun sin impacto economico sistemico -> 0

SI es riesgo economico (y debe puntuar ALTO):
- Petroperu: rescate, inyeccion fiscal, quiebra, deuda, crisis -> 75-90
- Camisea / gas natural: interrupcion, escasez, corte de suministro -> 75-90
- Conflictos mineros (Las Bambas, Tia Maria, Cuajone, Southern,
  Quellaveco) -> 55-70
- Paro nacional / huelga general -> 60-70
- Bloqueo de carreteras / interrupcion de cadena de suministro -> 40-55
- Consejo Fiscal advierte sobre gasto excesivo -> 50-60
- Gasto publico excesivo aprobado por Congreso / MEF -> 55-65
- Escasez de productos basicos / desabastecimiento -> 60-70
- Inflacion acelerando / BCRP subiendo tasa -> 45-55
- Tipo de cambio bajo presion / sol depreciandose fuerte -> 45-55
- Fuga de capitales / salida de inversion extranjera -> 65-75
- Inversion cancelada por conflicto social -> 50-60

Escala:
0 = Sin riesgo economico O noticia economica positiva/rutinaria
1-20 = Senal debil (volatilidad normal, advertencias menores)
21-40 = Estres moderado (indicadores deteriorandose, alzas en canasta
        basica)
41-60 = Vulnerabilidad significativa (disrupciones sectoriales, bloqueos)
61-80 = Crisis severa (colapso sectorial, paralisis energetica, rescate
        fiscal)
81-100 = Emergencia (default soberano, colapso financiero, hiperinflacion)

REGLA: Riesgo economico = DISRUPCION o AMENAZA ACTIVA a la actividad
economica. NO es lo mismo que "noticia sobre economia". Si no hay
disrupcion, amenaza o shock, el puntaje es 0.
\end{lstlisting}

\subsection*{User Message Template}
\label{app:user_template}

The user message is formatted as follows for a batch of $n$ articles:

\begin{lstlisting}
Evalua los siguientes {n} articulos. Para CADA articulo, responde con
un JSON en una linea separada. EXACTAMENTE {n} lineas, una por articulo,
en orden.

Articulo 1:
Titulo: {title_1}
Fuente: {source_1}
{summary_1[:300]}

Articulo 2:
...

Responde EXACTAMENTE {n} lineas JSON, una por articulo, en orden:
{"score": <0-100>}
{"score": <0-100>}
...
\end{lstlisting}

% ──────────────────────────────────────────────────────────────────────────────
\section{Post-Filter Rules}
\label{app:postfilter}

This appendix summarizes the deterministic post-filter rules applied after
AI classification.  All rules are implemented in the
\texttt{post\_filter\_scores()} function in
\texttt{src/nlp/classifier.py}.  Rules are applied in the order listed.
A ``crisis guard'' pattern prevents over-zeroing: articles containing
terms such as \textit{crisis, emergencia, bloqueo, paro, huelga,
desabastecimiento, escasez} are exempt from most zeroing rules.

\subsection*{Zeroing Rules (set both scores to 0)}

\begin{enumerate}
  \item \textbf{Sports}: Regex covering Liga~1 clubs,
    Copa Libertadores, UEFA, NBA, NFL, MLB, F1, Olympics, major athletes
    (Messi, Ronaldo, etc.) and match vocabulary (tabla de posiciones,
    jornada, fase de grupos).

  \item \textbf{Celebrity/tabloid}: Personal life news, entertainment
    shows (Esto es Guerra, Gran Hermano), music celebrities (Taylor Swift,
    Bad Bunny), and street crime (asaltos, accidentes de transito).

  \item \textbf{Foreign leader health}: Health/medical news about
    ex-presidents from Brazil, Argentina, Colombia, etc.\ (unless the
    title contains crisis vocabulary).

  \item \textbf{Cultural/arts events}: Theater festivals, ballet
    premieres, gallery openings, FAE Lima.

  \item \textbf{Intellectual/cultural deaths}: Obituaries of
    philosophers, writers, poets, musicians, and other cultural figures.

  \item \textbf{Foreign weather/local news}: Weather alerts in US
    cities, Argentine or Chilean local disasters.

  \item \textbf{Foreign domestic politics}: Elections or legislative
    debates in Colombia, Argentina, Brazil, USA, UK, EU (crisis exception
    preserved for tariff or trade news explicitly mentioning Peru).

  \item \textbf{Entertainment extras}: Concerts in Lima, new album
    releases, Oscar/Grammy ceremonies, box office records.

  \item \textbf{Individual petty crime}: Single muggings, bus thefts
    (distinct from organized crime, extortion, and crime rates which
    retain scores).
\end{enumerate}

\subsection*{Economic Score Zeroing Rules}

\begin{enumerate}
  \item \textbf{Foreign stock markets}: Wall Street, Dow Jones, Nikkei,
    Asian/European market summaries.
  \item \textbf{Foreign economic statistics}: US/EU/China unemployment,
    GDP, inflation.
  \item \textbf{Daily market summary columns}: ``Mercados e indicadores''
    routine reports.
  \item \textbf{Corporate positive growth}: Record sales, earnings growth,
    revenue increases.
  \item \textbf{Consumer/regulatory information}: SBS rule publications,
    Indecopi consumer sanctions, Sunafil guidance.
  \item \textbf{Health information}: Routine medical awareness articles.
  \item \textbf{Foreign IMF/World Bank loans to non-Peru countries}.
  \item \textbf{Vehicle recalls/product withdrawals}.
  \item \textbf{Gender diversity/talent features}.
  \item \textbf{Marketing/promotional content}: Sales, discounts,
    store openings.
  \item \textbf{Elections administrative procedures}: ONPE logistics,
    JNE candidate registration, RENIEC.
  \item \textbf{BCRP routine communications}: New banknote designs,
    reserve statistics.
  \item \textbf{Commercial chain openings}: Restaurant/gym expansions.
  \item \textbf{Economist opinion cap}: Named economist opinion columns
    capped at eco~=~20.
\end{enumerate}

\subsection*{Score Boost Rules (minimum floors)}

\begin{enumerate}
  \item \textbf{Political inversion correction}: Judicial/political
    confrontation articles with eco~$>$~0 and pol~$<$~65: set eco~=~0,
    pol~$\geq$~65.
  \item \textbf{Voter guide correction}: Ballot/voting guide articles:
    eco~=~0, pol~$\geq$~60.
  \item \textbf{Presidential vacancy}: Articles mentioning mocion de
    vacancia: pol~$\geq$~90, eco~$\geq$~40.
  \item \textbf{Presidential censure/removal}: pol~$\geq$~85, eco~$\geq$~40.
  \item \textbf{Premier removal/censure}: pol~$\geq$~65.
  \item \textbf{Minister removal/interpelacion}: pol~$\geq$~55.
  \item \textbf{Interpelacion + energy/economic crisis}: eco~$\geq$~55.
  \item \textbf{Camisea/gas crisis}: Corte, restriccion, emergencia
    related to Camisea or GNV: eco~$\geq$~80.
  \item \textbf{Petroperu crisis}: Rescate, inyeccion, quiebra,
    crisis related to Petroperu: eco~$\geq$~75.
\end{enumerate}

% ──────────────────────────────────────────────────────────────────────────────
\section{Extended Parameter Sensitivity}
\label{app:sensitivity_extended}

Table~\ref{tab:sensitivity_full} reports the full parameter sensitivity
grid for the week of March~14--18,~2026, covering all 27 combinations of
$\beta_{\text{pol}} \in \{0.3, 0.5, 0.7\}$,
$\beta_{\text{eco}} \in \{0.2, 0.3, 0.5\}$, and
$W \in \{60, 90, 120\}$.

\begin{table}[h!]
\centering
\caption{Full Parameter Sensitivity Grid: Week-Average IRP and IRE, March 14--18, 2026}
\label{tab:sensitivity_full}
\small
\begin{tabular}{ccccc}
\toprule
$\beta_{\text{pol}}$ & $\beta_{\text{eco}}$ & Window & Avg.\ IRP & Avg.\ IRE \\
\midrule
0.3 & 0.2 &  60 & 148.0 & 128.8 \\
0.3 & 0.2 &  90 & 183.8 & 163.2 \\
0.3 & 0.2 & 120 & 197.1 & 194.6 \\
0.3 & 0.3 &  60 & 148.0 &  94.9 \\
0.3 & 0.3 &  90 & 183.8 & 120.2 \\
0.3 & 0.3 & 120 & 197.1 & 143.4 \\
0.3 & 0.5 &  60 & 148.0 &  51.6 \\
0.3 & 0.5 &  90 & 183.8 &  65.3 \\
0.3 & 0.5 & 120 & 197.1 &  77.9 \\
\textbf{0.5} & 0.2 &  60 & 104.0 & 128.8 \\
\textbf{0.5} & 0.2 &  90 & 129.2 & 163.2 \\
\textbf{0.5} & 0.2 & 120 & 138.5 & 194.6 \\
\textbf{0.5} & \textbf{0.3} &  60 & 104.0 &  94.9 \\
\textbf{0.5} & \textbf{0.3} & \textbf{90}  & \textbf{129.2} & \textbf{120.2} \\
\textbf{0.5} & \textbf{0.3} & 120 & 138.5 & 143.4 \\
\textbf{0.5} & 0.5 &  60 & 104.0 &  51.6 \\
\textbf{0.5} & 0.5 &  90 & 129.2 &  65.3 \\
\textbf{0.5} & 0.5 & 120 & 138.5 &  77.9 \\
0.7 & 0.2 &  60 &  73.2 & 128.8 \\
0.7 & 0.2 &  90 &  91.0 & 163.2 \\
0.7 & 0.2 & 120 &  97.6 & 194.6 \\
0.7 & 0.3 &  60 &  73.2 &  94.9 \\
0.7 & 0.3 &  90 &  91.0 & 120.2 \\
0.7 & 0.3 & 120 &  97.6 & 143.4 \\
0.7 & 0.5 &  60 &  73.2 &  51.6 \\
0.7 & 0.5 &  90 &  91.0 &  65.3 \\
0.7 & 0.5 & 120 &  97.6 &  77.9 \\
\bottomrule
\multicolumn{5}{l}{\footnotesize Bold row is the baseline specification.}
\end{tabular}
\end{table}

% ──────────────────────────────────────────────────────────────────────────────
\section{Source-Level Monthly Statistics}
\label{app:source_monthly}

Table~\ref{tab:sources_monthly} provides monthly IRP and IRE averages for
the full 443-day sample.

\begin{table}[h!]
\centering
\caption{Monthly Average IRP, IRE, and Article Volume}
\label{tab:sources_monthly}
\begin{tabular}{lrrr}
\toprule
Month & Avg.\ IRP & Avg.\ IRE & Mean daily articles \\
\midrule
2025-01 & 73.2 &  61.3 &  50.8 \\
2025-02 & 50.8 &  33.4 &  38.6 \\
2025-03 & 86.9 &  38.2 &  43.0 \\
2025-04 & 62.9 &  47.0 &  35.4 \\
2025-05 & 64.1 &  45.8 &  30.4 \\
2025-06 & 89.7 &  26.5 &  31.8 \\
2025-07 & 77.5 &  47.4 &  39.4 \\
2025-08 & 81.3 &  31.7 &  39.5 \\
2025-09 & 65.5 &  90.1 &  47.7 \\
2025-10 & 76.2 &  82.7 &  61.2 \\
2025-11 & 61.0 &  48.5 &  62.0 \\
2025-12 & 50.3 &  36.7 &  67.3 \\
2026-01 & 67.5 &  61.1 &  79.9 \\
2026-02 & 97.5 &  60.7 &  98.7 \\
2026-03 & 85.6 & 213.7 & 337.6 \\
\bottomrule
\multicolumn{4}{l}{\footnotesize March 2026 reflects elevated article volume from multiple concurrent crises.}
\end{tabular}
\end{table}

The dramatic increase in mean daily articles from 30--50 (early 2025) to
337.6 (March 2026) reflects expanded pipeline coverage: by early 2026,
all 16 sources were contributing daily, while in early 2025 only a subset
were active (see Table~\ref{tab:sources}, column ``Days active'').  The
March 2026 IRE average of 213.7 is nearly three times any prior month,
driven by the Camisea gas crisis.

\end{document}
